\documentclass[11pt]{article}
\usepackage[T1]{fontenc}
\usepackage{amsmath,amssymb,amsthm,hyperref}
\newtheorem{conjecture}{Conjecture}
\newtheorem{lemma}{Lemma}
\newtheorem{proposition}{Proposition}
\newtheorem{corollary}{Corollary}
\title{Strong relative cone domains: an intake contract and elementary reductions}
\author{Research working draft; originality unassessed}
\date{November 22, 2026}
\begin{document}
\maketitle
\begin{abstract}
This is an initial research note, not a proof of the strong relative Morrison--Kawamata movable cone conjecture in relative dimension at most two. We fix the requested effective cone, prove a finite-stabilizer refinement lemma and an elementary pointed-cone criterion, and exhibit an abstract obstruction to saturated lifting. The geometric hypotheses needed to apply these reductions to all target fibrations remain unproved. The accompanying verification verdict is FAIL.
\end{abstract}
% Original target, deliberately stated as an unproved problem.
\section{The fixed target}\label{sec:target}
\begin{conjecture}[Research target; not proved here]\label{conj:target}
Let $f:(X,\Delta)\to S$ be a complex terminal Calabi--Yau fibre space in the sense specified in the problem contract, with $X$ normal, projective and $\mathbb Q$-factorial, $K_X+\Delta\equiv_S0$ and $\dim X-\dim S\leq2$.
Put $V=N^1(X/S)_{\mathbb R}$, $C_e=\overline{\operatorname{Mov}}(X/S)\cap\operatorname{Eff}(X/S)$ and let $\Gamma\subset\operatorname{GL}(V)$ be the image of $\operatorname{PsAut}(X/S,\Delta)$.
There is a rational polyhedral cone $P\subset C_e$ such that $\Gamma P=C_e$, $\gamma P\cap\operatorname{Int}(P)=\varnothing$ for every $\gamma$ not stabilizing $P$, and $\operatorname{Stab}_\Gamma(P)=\{1\}$. No nondegeneracy or $\mathbb Q$-factorial-base hypothesis is imposed.
\end{conjecture}
\noindent The extra convention checks in the fixed problem are obligations, not permission to restrict the class. Conjecture~\ref{conj:target} remains open in this note.

\section{Literature boundary}
Li--Zhao distinguish weak and strong domains and analyze degenerate-cone quotients \cite{LZ}. Li's sequel supplies generic-to-relative machinery with a modified-cone conclusion \cite{Li}; Moraga--Stark's geometric statement concerns chamber and face orbits \cite{MS}. Faucher again separates weak and nondegenerate strong conclusions \cite{F}. These statements must be compared on exact hypotheses and cones. The publisher texts and their complete version differences have not yet been independently checked in this archive. In particular, no claim of complete novelty certification is made. Kawamata's examples remain essential input for boundary checks \cite{K}.
\section{Abstract reductions with complete proofs}
All interiors below are ambient interiors. These lemmas do not assert that their hypotheses hold for the geometric target. No originality is claimed for these elementary reductions.

\begin{lemma}[Finite-stabilizer refinement]\label{lem:finite}
Let $V_{\mathbb Q}$ be finite-dimensional, let $\Gamma\subset\operatorname{GL}(V_{\mathbb Q})$, and let $C\subset V_{\mathbb R}$ be invariant and have nonempty interior. Suppose a full-dimensional rational polyhedral cone $P\subset C$ satisfies $\Gamma P=C$ and, for every $g\in\Gamma$, either $gP=P$ or $gP\cap\operatorname{Int}(P)=\varnothing$. If $H=\operatorname{Stab}_\Gamma(P)$ is finite, then $C$ has a strong rational polyhedral fundamental domain for $\Gamma$.
\end{lemma}
\begin{proof}
For each $h\in H\setminus\{1\}$, the rational subspace of covectors fixed by $h$ is proper, because the action on $V$ is faithful. A finite union of proper rational subspaces cannot exhaust $V_{\mathbb Q}^*$. Choose a rational covector $\ell$ outside their union and put
\[
 D=P\cap\bigcap_{h\in H}\{x:\ell(x)\leq\ell(hx)\}.
\]
This is rational polyhedral. For every $x\in P$, choose $h_0$ minimizing $\ell(hx)$ among $h\in H$. Then $h_0x\in D$, so $P=HD$ and $\Gamma D=C$. The cone $D$ is full-dimensional: otherwise the finite union $HD$ lies in a finite union of proper linear subspaces and cannot cover $\operatorname{Int}(P)$.

For $h\in H\setminus\{1\}$ and $y\in\operatorname{Int}(D)$ one has $\ell(y)<\ell(hy)$. Indeed, equality for a nonzero linear form $\ell\circ h-\ell$ cannot occur at an ambient interior point of a set contained in its nonnegative half-space. If $y\in hD\cap\operatorname{Int}(D)$ and $x=h^{-1}y\in D$, the strict inequality for $h^{-1}$ gives $\ell(y)<\ell(x)$, whereas $x\in D$ gives $\ell(x)\leq\ell(hx)=\ell(y)$, a contradiction. For $g\notin H$, $gD\subset gP$ and $\operatorname{Int}(D)\subset\operatorname{Int}(P)$, so the assumed weak separation gives $gD\cap\operatorname{Int}(D)=\varnothing$. Thus this exclusion holds for every nonidentity $g$; since $D$ has nonempty interior, its stabilizer is trivial.
\end{proof}

\begin{lemma}[Pointed polyhedral stabilizer]\label{lem:pointed}
Let a lattice $\Lambda\subset V$ be preserved by $\Gamma\subset\operatorname{GL}(V)$. If $P$ is full-dimensional, pointed and rational polyhedral with respect to $\Lambda\otimes_{\mathbb Z}\mathbb Q$, then $\operatorname{Stab}_\Gamma(P)$ is finite.
\end{lemma}
\begin{proof}
A stabilizer element permutes the finitely many extremal rays. On each ray there is a unique primitive lattice generator. An invertible lattice map sends it to the primitive generator of the image ray. The induced permutation action is faithful, since the ray generators span $V$. The stabilizer therefore embeds in a finite symmetric group. In dimension zero the image group is trivial and the assertion holds separately.
\end{proof}

\begin{corollary}\label{cor:pointed}
If $C\cap(-C)=0$ and a lattice-preserving action on $C$ admits a full-dimensional weak rational polyhedral domain, that domain can be refined to a strong one, even when the closure of $C$ has nonzero lineality.
\end{corollary}
\begin{proof}
Every subcone of $C$ is pointed. Apply Lemmas~\ref{lem:pointed} and \ref{lem:finite}.
\end{proof}

\begin{proposition}[An obstruction to a universal saturated-lift argument]\label{prop:shear}
Let $V=\mathbb R^2$, $C=\{(x,y):y\geq0\}$ and $g(x,y)=(x+y,y)$. The action of $\Gamma=\langle g\rangle\subset\operatorname{GL}_2(\mathbb Z)$ has a strong rational polyhedral domain for the quotient action image on the quotient by $W=\mathbb R(1,0)$, but it has no convex conical strong domain on $C$. This is an abstract example, not a Calabi--Yau counterexample.
\end{proposition}
\begin{proof}
The quotient action image is trivial and the quotient cone is $\mathbb R_{\geq0}$, its own domain. Suppose $P\subset C$ is a convex cone with $\Gamma P=C$. Every point of $W$ is fixed by every group element, hence coverage forces $W\subset P$. Coverage of $(0,1)$ also forces $P$ to contain some $v$ with positive second coordinate. Since $W+\mathbb R_{\geq0}v=C$, one gets $P=C$. Every $g^n$ stabilizes $P$, and $g\neq1$, contradicting the strong-domain requirement.
\end{proof}

\begin{proposition}[Why the effective boundary matters]\label{prop:open}
For the same action, $C_0=\{(x,y):y>0\}\cup\{0\}$ has the strong rational polyhedral domain $D=\{(x,y):0\leq x\leq y\}$. Here $C_0\cap(-C_0)=0$ but the closure has lineality $W$.
\end{proposition}
\begin{proof}
For $y>0$, subtracting $\lfloor x/y\rfloor y$ from $x$ puts $(x,y)$ in $D$, and $0\in D$. On interiors, $g^nD$ has $ny<x<(n+1)y$, which is disjoint from $0<x<y$ for $n\neq0$. More strongly, the closed intervals $[ny,(n+1)y]$ do not meet $(0,y)$ for nonzero integral $n$. Thus the exact strong separation holds. The assertions on lineality follow directly from the definitions.
\end{proof}

\section{The still missing geometric step}
The lemmas prove a sufficient conditional route: a weak rational polyhedral domain on the actual $C_e$ with finite stabilizer implies the original strong conclusion. Neither existence on this exact cone nor finite stabilizers for every target fibration has been proved here. Results for $\operatorname{Mov}_{+}$ cannot be substituted without proving the required comparison. No proof of Conjecture~\ref{conj:target} is supplied.

\section{Unresolved scope and verification}
The original target permits non-$\mathbb Q$-factorial bases, arbitrary absolute dimension and every admissible boundary. Neither restrictions to K3 fibres nor vanishing relative irregularity may be imposed as substitute targets. The exact terminal and linear-versus-numerical conventions must be reconciled from the originals. None of the elementary results above supplies that geometric reconciliation. The eight original gates remain unsatisfied as a conjunction.
\begin{thebibliography}{9}
\bibitem{LZ} Z. Li and H. Zhao, \emph{On the relative Morrison--Kawamata cone conjecture}, Proc. London Math. Soc. 131 (2025), e70099. doi:10.1112/plms.70099. arXiv:2206.13701v8.
\bibitem{Li} Z. Li, \emph{On the relative Morrison--Kawamata cone conjecture (II)}, arXiv:2309.04673v1. Journal metadata: J. Algebraic Geom. 36 (2027), 1--50, online August 24, 2026, doi:10.1090/jag/864. Publisher body not independently compared here.
\bibitem{MS} J. Moraga and T. Stark, \emph{The geometric cone conjecture in relative dimension two}, arXiv:2409.13068v1 (2024).
\bibitem{F} A. Faucher, \emph{The relative movable cone conjecture for K-trivial fibrations in varieties with well-clipped movable cones}, arXiv:2603.22457v1 (2026).
\bibitem{K} Y. Kawamata, \emph{On the cone of divisors of Calabi--Yau fiber spaces}, Internat. J. Math. 8 (1997), 665--687, arXiv:alg-geom/9701006.
\end{thebibliography}
\end{document}
